\begin{abstract}
Predicting dough deformation during robotic manipulation requires a model that captures recoverable, rate-dependent, and persistent deformation without making calibration impractical. This paper presents TaichiDough, a framework for differentiable identification of effective simulation parameters from single-view depth observations and recorded rigid-tool trajectories. A table-aligned reconstruction initializes a three-dimensional moving least squares material point simulator with affine particle--grid transfer. The constitutive approximation combines fixed-corotated elasticity, an additive viscous term, and principal-stretch plastic projection. Scene geometry and marker-to-tool transforms are calibrated separately and held fixed during material identification. A partial-observation objective combines depth-change agreement, observed-region coverage, and directed distances to the predicted visible geometry, without treating unobserved space as known empty space. Checkpointed reverse-mode differentiation supports bounded optimization of Young's modulus, Poisson's ratio, viscosity, and two plastic stretch limits, with an extension to shared-material fitting across recordings. The central research question is whether time-resolved visual calibration improves prediction of unseen dough manipulation under fixed parameters. We distinguish this predictive objective from recovery of unique physical constants, since hidden volume, contact assumptions, and incomplete mechanical excitation can produce similar observations for different parameter combinations. The framework defines an evaluation protocol for testing that distinction; real-data predictive performance remains to be established.
\end{abstract}
